\documentclass[11pt,a4paper]{article}
\usepackage[utf8]{inputenc}
\usepackage{graphicx}
\usepackage[left=2.5cm,top=3cm,right=2.5cm,bottom=3cm,bindingoffset=0.5cm]{geometry}
\usepackage{AEDLogica, AEDEspecificacion, AEDTADs}
\usepackage{caratula}
\newlength{\miSangria}
\setlength{\miSangria}{2em}

% Asi pueden escribir nuevos comandos. 
% Este por ejemplo asegura q los nombres 
% que figuren con una tipografia diferenciada  
\newcommand{\Tipo}[1]{\mathsf{#1}}
% la sintaxis es \newcommand{\nombreDeLaMacro}[cantidadDeParametros]{Lo que va ser remplazado por el macro} 
\newcommand{\norm}[1]{\vert#1\vert}


\begin{document}

\section{ElMasCanjeador}
	\vspace{1.0cm}

	\begin{proc}{ElMasCanjeador}
		{
			\In AM: \Tipo{Airmiles}
		}
		{\Tipo{\Z}}
		\requiere{ 
            \norm{AM.Usuarios} > 0 \land \norm{AM.Historial} > 0
        }
		\aseguraLargo{
            \existe{id}{\Tipo{\Z}}{id \in claves(AM.Usuarios) \yLuego res = id
            \\ 
            \land \paraTodo{id2}{\Tipo{\Z}}{id2 \in claves(AM.Usuarios) \implicaLuego
            \\
            MillasPerdidas(AM.Historial,res,AM.Usuarios) \ge
            \\
            MillasPerdidas(AM.Historial,id2,AM.Usuarios)}}
		}
	\end{proc}

    \comentario{Es similar a EsTransferencia pero aca no importa el destino, unicamente chequeo que exista el usuario de origen.}
    \\
    \predLargo{EsTransferenciaEnviada}
    {Transaccion: \Tipo{Transacciones}, idOrigen: \Tipo{\Z}, Usuarios: \Tipo{\dict{\Z}{Categoria}}}
    {
        Transaccion.UsuarioOrigen = idOrigen \land
        \\
        Transaccion.TipoDeTransaccion = 1 \land
        \\
        Transaccion.UsuarioDestino \in claves(Usuarios) \land
        \\
        Transaccion.UsuarioDestino \neq idOrigen
        \\
        \comentario{Deberia verificar los otros campos? Se complica...}
    }

    \vspace{1cm}
    \auxLargo{MillasPerdidas}
    {Historial: \Tipo{\seq{Transacciones},idUsuario: \Tipo{\Z},Usuarios: \Diccionario{\Tipo{\Z}}{categoria}}}
    {\Tipo{\R}}
    {\sum\limits_{i=0}^{\norm{Historial}-1} 
    \IfThenElse
    {EsTransferenciaEnviada(Historial[i],idUsuario,Usuarios) \lor 
    \\
    EsCanjeo(Historial[i],idUsuario)}
    {Historial[i].Millas}
    {0}
    }

\end{document}